\chapter{Literature Review}

\section{The Malaysian Context}
The literature establishes that while global ASD prevalence rises, the Malaysian landscape presents unique logistical and cultural barriers \citep{Lim2021_79, Ng2021_114}. The lack of localized support infrastructures necessitates the exploration of virtual modalities that account for regional digital divides.

\section{Virtual Counseling Interventions}
Foundational telehealth frameworks are heavily Western-centric. A review of recent literature indicates that while telehealth mitigates caregiver burden, current implementations suffer from methodological myopia by relying on generalized statistical regressions that assume caregiver uniformity \citep{Malgaroli2023_92, Kolog2018_108}.

\section{The Autism Caregiver and the TPB}
To accurately model the decision-making process of caregivers, the Theory of Planned Behavior \citep{ajzen1991} provides a robust framework. The constructs of Attitude, Subjective Norms, and Perceived Behavioral Control (PBC) serve as the primary variables dictating caregiver engagement with virtual support.

\section{Clustering Methods in Social Sciences}
Despite the robust application of cluster analysis in family typologies, its application to caregiver psychology remains nascent. By mathematically transforming ordinal survey data via Partial Least Squares Structural Equation Modeling (PLS-SEM) into continuous latent variables, this study sets the foundation for advanced segmentation \citep{Lanza2003_71}. Applying a comparative framework of deterministic and probabilistic clustering algorithms—specifically K-Means, Latent Profile Analysis (LPA), and Two-Stage Clustering—to segment caregivers based on TPB constructs represents a highly novel methodological contribution.
